% sec:ad-christoffels — AD for Christoffel Symbols
\section{AD for Christoffel Symbols}
\label{sec:ad-christoffels}

At every integration step, the geodesic right-hand side needs the
40 metric derivatives $\pd{\mu} g_{\alpha\beta}$ --- and the
dual-number machinery of the previous section produces them
automatically.  Two routes connect this output to equations of
motion: one that assembles Christoffel symbols explicitly, and
one that --- by exploiting the covariant Hamiltonian formulation
of \cref{sec:geo-formulations} --- bypasses them entirely.  Both
start from the same four seeded evaluations of the metric
function, costing roughly eight metric evaluations in total
(\cref{sec:ad-forward}); they differ in how the resulting
derivatives are consumed.  The second route is what the codebase
uses, and it reveals that the right choice of formulation can
eliminate an entire layer of intermediate computation.

\paragraph{Assembling Christoffel symbols.}

The Christoffel symbols defined in \cref{sec:dg-covariant}
(\cref{eq:christoffel-formula}) are
%
\begin{equation}
  \christoffel{\mu}{\alpha}{\beta}
    = \tfrac{1}{2}\,g^{\mu\nu}\bigl(
        \pd{\alpha} g_{\nu\beta}
      + \pd{\beta} g_{\nu\alpha}
      - \pd{\nu} g_{\alpha\beta}
    \bigr) \,.
  \label{eq:christoffel-recall}
\end{equation}
%
Each component is a contraction of metric derivatives against the
inverse metric, summed over $\nu = 0, \ldots, 3$.  The required
inputs are the four symmetric derivative tensors
$\pd{0} g_{\alpha\beta}, \ldots, \pd{3} g_{\alpha\beta}$ --- 40
independent values from four seeded evaluations
(\cref{eq:dual-seed}) --- plus the inverse metric $\inversemetric$,
computed separately since matrix inversion is not part of the AD
pipeline.

The function \texttt{christoffelsAD} in
\codemodule{Spacetime.AutoDiff} implements
\cref{eq:christoffel-recall} directly:
\coderef{Spacetime.AutoDiff}{christoffelsAD}
%
\begin{lstlisting}[language=Haskell]
christoffelsAD
  :: (forall a. Floating a => [a] -> [a])
  -> Sym4 'Contravariant -> [Double]
  -> Christoffel4
christoffelsAD metricFn ginv coords =
  let (dg0, dg1, dg2, dg3) =
        metricDerivativesAD metricFn coords
      dgs = [dg0, dg1, dg2, dg3]
      gamma mu a b =
        0.5 * sum [ sym4Get ginv mu rho
          * ( sym4Get (dgs !! a) rho b
            + sym4Get (dgs !! b) rho a
            - sym4Get (dgs !! rho) a b )
          | rho <- [0..3] ]
  in Christoffel4 (buildSym4 0)
       (buildSym4 1) (buildSym4 2)
       (buildSym4 3)
\end{lstlisting}
%
The list comprehension in \texttt{gamma} is a direct transcription
of \cref{eq:christoffel-recall}: the three \texttt{sym4Get} terms
inside the parentheses correspond to $\pd{\alpha} g_{\nu\beta}$,
$\pd{\beta} g_{\nu\alpha}$, and $-\pd{\nu} g_{\alpha\beta}$
respectively, contracted against $g^{\mu\nu}$ via the outer sum
over $\nu$ (here called \texttt{rho}).  The result is a
\texttt{Christoffel4} value: four \texttt{Sym4} tensors, one per
upper index~$\mu$, storing 40 independent components.
\implnote{The symmetry
$\christoffel{\mu}{\alpha}{\beta}
= \christoffel{\mu}{\beta}{\alpha}$, inherited from the symmetry
of $g_{\alpha\beta}$, lets \texttt{Christoffel4} store
40 rather than 64 components --- one symmetric $4 \times 4$
matrix per upper index.}

\paragraph{The covariant shortcut.}

Christoffel symbols are the natural intermediary when the equations
of motion take the second-order form
$\ddot{x}^\mu
= -\christoffel{\mu}{\alpha}{\beta}\,
  \dot{x}^\alpha\,\dot{x}^\beta$,
but the covariant momentum equation derived in
\cref{sec:geo-formulations} provides a shorter path:
%
\begin{equation}
  \dot{p}_\mu
    = \tfrac{1}{2}\,
      (\pd{\mu} g_{\alpha\beta})\,
      \dot{x}^\alpha\,\dot{x}^\beta \,.
  \label{eq:pdot-ad}
\end{equation}
%
This equation uses only the raw metric derivatives
$\pd{\mu} g_{\alpha\beta}$ --- exactly what
\texttt{metricDerivativesAD} returns --- without the additional
contraction against $\inversemetric$ that Christoffel assembly
requires.  The covariant momentum $p_\mu$ already carries a
lowered index, so the equation of motion never needs to raise
indices with $g^{\mu\nu}$ --- precisely the contraction that
Christoffel symbols encode.  The velocity
$\dot{x}^\mu = g^{\mu\nu}\,p_\nu$ still needs the inverse
metric, but that is a single matrix-vector product, far simpler
than assembling and then contracting 40 Christoffel components.
\physnote{\Cref{eq:pdot-ad} restates the covariant momentum
equation (\cref{eq:geo-pdot-cov}).  It follows directly from
Hamilton's equations without introducing Christoffel symbols:
the lowered momentum absorbs the connection terms that would
otherwise appear in the second-order form.}

The codebase exploits this structure.  The geodesic
integrator never calls \texttt{christoffelsAD}.  Instead,
the right-hand-side function in
\codemodule{Geodesic.Hamiltonian} calls
\texttt{metricDerivativesAD} directly, computes
$\dot{x}^\mu = g^{\mu\nu}\,p_\nu$, and evaluates
\cref{eq:pdot-ad} by contracting each derivative tensor
against the velocity.
\coderef{Geodesic.Hamiltonian}{hamiltonianRHS}
No Christoffel symbols are assembled at any stage --- the raw
metric derivatives flow straight into the momentum update.

\paragraph{Partial application of physical parameters.}

Each spacetime's \texttt{Spacetime} instance uses the
partial-application pattern from the previous section.
For Schwarzschild in Kerr--Schild coordinates:
%
\begin{lstlisting}[language=Haskell]
christoffel4 sks pos =
  let metricFn :: forall a. Floating a
               => [a] -> [a]
      metricFn = schwarzschildMetricFn
                   (realToFrac m)
  in christoffelsAD metricFn ginv coords
\end{lstlisting}
%
The call \texttt{realToFrac} lifts the mass $M$ into the AD
library's numeric type; partial application then produces a
closure that treats $M$ as a constant during differentiation.
The same pattern applies to every spacetime in the codebase:
Kerr partially applies $M$ and $a$; Morris--Thorne partially
applies $b_0$.

\paragraph{The metric-function pipeline.}

The combination of AD and the covariant momentum equation produces
a clean separation of concerns: a new spacetime enters the ray
tracer as a single function --- ten numbers from four
coordinates --- and the entire integration pipeline follows
mechanically.  The AD machinery extracts derivatives; the
covariant formulation turns them into forces; the adaptive
Dormand--Prince integrator of \cref{sec:rk-dormand-prince}
advances the state.  Nothing in this chain depends on which
spacetime is being traced.  The effort of supporting a new spacetime scales
with the complexity of the metric itself, not with the complexity
of its 40 derivative expressions.  A finite-difference fallback
(\texttt{christoffelsFiniteDiff}) exists for spacetimes that lack
a polymorphic metric function, at the cost of truncation error
discussed in \cref{sec:ad-comparison}.  The
\texttt{christoffelsAD} function remains available for tasks that
require assembled Christoffel symbols (computing curvature
invariants, verifying geodesic deviation), but the integrator's
hot path never uses it.
